% Lösungen Einheit 4 von 4 – Sachgeraden (zusammenbau v0.1)
\einheitenkopf{Einheit 4 von 4 · Sachgeraden}

\erg{18}{a) Strecke – das Seil beginnt am ersten und endet am zweiten Mast \quad b) $\vec x = (0 | 0 | 8) + t \cdot (6 | 8 | -3)$ mit $0 \le t \le 1$ \quad c) Unteres Seil: $\vec x = (0 | 0 | 6) + s \cdot (4 | 8 | -2)$. Gleiche zweite Koordinate: $8s = 2$ liefert $s = 0{,}25$; gleiche erste: $-2 + 3t = 1$ liefert $t = 1$; Punkte $(1 | 2 | 5{,}5)$ und $(1 | 2 | 10)$, Abstand $10 - 5{,}5 = 4{,}5$ m \quad d) Horizontalabstand $\sqrt{36 + 64} = 10$ m, Höhenunterschied $1{,}5$ m, Neigung $\frac{1{,}5}{10} = 15\,\%$ \quad e) Aus dem bekannten Teilverhältnis den Auflagepunkt $X$ auf $UV$ bestimmen; die Seilgerade durch den Befestigungspunkt am ersten Mast und $X$ aufstellen; den zweiten Mast als senkrechte Gerade durch $M$ beschreiben (erste und zweite Koordinate wie bei $M$) und mit der Seilgeraden schneiden; der gesuchte Abstand ist die dritte Koordinate des Schnittpunkts minus die Dachhöhe – eine Höhendifferenz, kein Lotabstand zu einer Ebene.}
\erg{19}{a) Horizontalabstand $\frac{0{,}9 - 0{,}3}{0{,}04} = 15$ m; linkes Ende bei $x = -15$, also $15 - 10 = 5$ m vom Rand entfernt}
\erg{20}{a) Gefragt ist der senkrechte Abstand: Punkte mit gleicher erster und zweiter Koordinate suchen und die dritten Koordinaten subtrahieren. Der Abstand windschiefer Geraden ist die kürzeste, meist schräge Verbindung – ein anderer Wert. \quad b) Senkrecht übereinander heißt: gleiche erste und gleiche zweite Koordinate – das sind zwei Gleichungen; die dritte Koordinate bleibt frei, ihre Differenz ist der Höhenunterschied. \quad c) $x_3 = 0{,}5 + 0{,}6t = 5$ liefert $t = 7{,}5$ – nach $7{,}5$ Minuten, also um 12.07 Uhr und 30 Sekunden.}
